\documentclass[12pt,a4paper]{article}

% Pacotes básicos
\usepackage[utf8]{inputenc}
\usepackage[T1]{fontenc}
\usepackage[brazil]{babel}
\usepackage{amsmath,amssymb,amsthm}
\usepackage{graphicx}
\usepackage{hyperref}
\usepackage{natbib}
\usepackage{geometry}
\usepackage{setspace}

% Configurações
\geometry{margin=2.5cm}
\onehalfspacing

% Informações do artigo
\title{\textbf{Título do Seu Paper}}
\author{Nome do Autor\thanks{Email: autor@email.com}}
\date{\today}

\begin{document}

\maketitle

\begin{abstract}
Resumo do seu paper. Descreva brevemente o objetivo, metodologia, 
resultados principais e conclusões. Mantenha entre 150-250 palavras.

\textbf{Palavras-chave:} palavra1, palavra2, palavra3
\end{abstract}

\section{Introdução}
\label{sec:intro}

Inclua a introdução aqui. Você pode usar \cite{autor2024} para citar referências.

% Introdução - Seção separada para facilitar edição

Este é um arquivo de exemplo para a seção de introdução.
Separe cada seção em arquivos diferentes para facilitar:
\begin{itemize}
    \item Edição colaborativa
    \item Revisão por IAs
    \item Controle de versão
\end{itemize}

A introdução deve apresentar:
\begin{enumerate}
    \item O problema de pesquisa
    \item A relevância do estudo
    \item Os objetivos
    \item A estrutura do paper
\end{enumerate}


\section{Metodologia}
\label{sec:metodo}

Descreva sua metodologia. Equações podem ser incluídas assim:

\begin{equation}
    E = mc^2
    \label{eq:einstein}
\end{equation}

% Metodologia - Seção separada para facilitar edição

Descreva detalhadamente:
\begin{itemize}
    \item Abordagem de pesquisa
    \item Dados utilizados
    \item Ferramentas e técnicas
    \item Procedimentos de análise
\end{itemize}

\subsection{Materiais}

Liste os materiais, equipamentos ou datasets utilizados.

\subsection{Procedimentos}

Descreva o passo a passo da metodologia.


\section{Resultados}
\label{sec:resultados}

Apresente seus resultados. Figuras podem ser incluídas:

\begin{figure}[h]
    \centering
    \includegraphics[width=0.8\textwidth]{figures/exemplo.png}
    \caption{Legenda da figura}
    \label{fig:exemplo}
\end{figure}

% Resultados - Seção separada para facilitar edição

Apresente os resultados de forma clara e objetiva.

\subsection{Análise Quantitativa}

Tabelas podem ser incluídas:

\begin{table}[h]
    \centering
    \caption{Exemplo de tabela}
    \begin{tabular}{|l|c|c|}
        \hline
        \textbf{Variável} & \textbf{Valor A} & \textbf{Valor B} \\
        \hline
        Item 1 & 10 & 20 \\
        Item 2 & 15 & 25 \\
        Item 3 & 20 & 30 \\
        \hline
    \end{tabular}
    \label{tab:exemplo}
\end{table}

\subsection{Análise Qualitativa}

Descreva os padrões e insights encontrados.


\section{Discussão}
\label{sec:discussao}

Discuta os resultados e suas implicações.

\section{Conclusão}
\label{sec:conclusao}

Apresente as conclusões principais e trabalhos futuros.

% Bibliografia
\bibliographystyle{apalike}
\bibliography{references}

\end{document}
